% ===================================================================
%  LENTELĖ Nr. 4 — Brandos amžius ir senatvė.
%  Traduction 2026 d'apres texto/io, controlee sur texto/fr, texto/ru
%  et texto/cs. Consigne : voir texto/lt/10-lentele-01.tex.
%
%  LA MALADIE ETAIT, POUR LE GALICIEN, L'UN DES ENDROITS OU LA
%  DISTANCE A LA LANGUE VOISINE ETAIT MINIMALE : on apprend le nom
%  des maux a l'ecole et chez le medecin, c'est-a-dire au meme
%  endroit qu'on apprend la langue d'a cote. LE LITUANIEN LE
%  CONFIRME PAR L'AUTRE BOUT. « reumatas », « pneumonija »,
%  « mikstūra », « receptas », « kalendorius », « portretas » entrent
%  du dehors — mais « karščiavimas » la fievre, « kosulys » la toux,
%  « svaigulys » le vertige, « sloga » le rhume, « dantų skausmas »
%  le mal de dents sont a lui. La ligne passe entre le mal QU'ON
%  SENT et le mal QU'ON DIAGNOSTIQUE : le premier a un nom depuis
%  toujours, le second arrive avec l'ordonnance.
%
%  ET C'EST LA MEME LIGNE QU'AU TABLEAU 2, ou le corps etait tout
%  lituanien et le sport tout emprunte. Le critere du lieu
%  d'apprentissage ne separe pas des domaines, il separe, DANS un
%  domaine, ce qu'on eprouve de ce qu'on vous apprend.
%
%  LE SAC DE CHASSE EN BANDOULIERE EST ENTIEREMENT EN GRAS, comme
%  dans les vingt colonnes qui l'ont corrige : « medžioklės krepšiu
%  per petį », le qualificatif fait corps avec le nom et le gras le
%  couvre en entier.
%
%  DEUX ORDRES QU'IL A FALLU TENIR SANS FORCER LA PHRASE. Le
%  sommelier debouche la bouteille (52) au tire-bouchon (51) et
%  l'aveugle (20) traverse la place devant la pharmacie (16), conduit
%  par son caniche (21) : dans les deux cas l'ordre lituanien le plus
%  spontane aurait interverti les deux renvois, et dans les deux cas
%  une autre construction, tout aussi lituanienne, garde l'ordre de
%  l'ido. On ne tord pas la langue pour effacer une inversion — voir
%  le tableau 3, ou l'on en a laisse deux — mais quand deux tours
%  valent l'un l'autre, on prend celui qui suit la source.
% ===================================================================

%%K t04-apar-1 apar
\VUsaut{4.0mm}
\VUtitre{15.0pt}{60}{LENTELĖ Nr. 4}
\VUsaut{1.6mm}
\VUtitre{11.4pt}{0}{Brandos amžius ir senatvė.}
\VUsaut{3.2mm}

%%K t04-tit-1 sub
\VUcentre{11.4pt}{0}{\textit{Pirmoji scena.}}
\VUsaut{1.6mm}
\VUtitre{11.4pt}{0}{Vestuvių pokylis.}
\VUsaut{2.6mm}

%%K t04-tit-2 sub
\VUpk{10.2pt}{40}{Ruduo.}
\VUsaut{3.0mm}

%%K t04-c1-01-1 p
1. --- Jaunuoliai užaugo. Vienas iš jų, Jonas, veda. Dalyvaujame \VUgras{vestuvių
pokylyje,} kuris aplink tą patį stalą surenka sutuoktinių šeimas.

%%K t04-c1-02-1 p
2. --- Esame \VUgras{rudenį,} \VUgras{rugsėjo mėnesį.} Šaltas \VUgras{spalio} ir
\VUgras{lapkričio} vėjas dar nenuplėšė nuo kaimo jo žalios lapijos. Vakaro oras
drungnas ir kvapnus; todėl vakarienė vyksta po \VUgras{veranda} \textsuperscript{(8)}
prie sodo.

%%K t04-c1-03-1 p
3. --- Toli, ant \VUgras{kalvos} \textsuperscript{(3)}, stovi Jono tėvų namai,
elegantiška \VUgras{pilis} \textsuperscript{(1)}, pasislėpusi žalumoje; gražūs
vynuogynai, duodantys puikaus vyno, dengia kalvelės šlaitą. Vynuogininkas juos
augina ir prižiūri.

%%K t04-c1-04-1 p
4. --- Virš vynuogynų praskrenda būrys \VUgras{kurapkų} \textsuperscript{(2)},
išgąsdintų besiartinančio \VUgras{eigulio} \textsuperscript{(5)}. Šis, su
\VUgras{šautuvu} po pažastimi ir \VUgras{medžioklės krepšiu per petį,} apeina
\VUgras{krūmyną} \textsuperscript{(4)} su savo \VUgras{šunimi} \textsuperscript{(6)}.
Jis prižiūri valdą ir neleidžia brakonieriams naikinti medžiojamųjų žvėrių.

%%K t04-c1-05-1 p
5. --- Už \VUgras{verandos} kopia \VUgras{vynmedžių špaleras,} nuo kurio kabo tankios
\VUgras{vynuogių kekės} \textsuperscript{(7)}.

%%K t04-c1-06-1 p
6. --- Gražios rudens gėlės, puošiančios \VUgras{stalą} \textsuperscript{(16)}, yra
\VUgras{chrizantemos} \textsuperscript{(9)}; jaunosios \VUgras{tėvas}
\textsuperscript{(10)} įsisegė vieną iš jų į savo frako atlapo kilpelę.

%%K t04-c1-07-1 p
7. --- Vaišės baigiasi. \VUgras{Tarnas} \textsuperscript{(11)} pradeda nešti deserto
lėkštes ir netrukus nurinks įvairias stalo įrankių dalis --- \VUgras{šaukštus}
\textsuperscript{(14)}, \VUgras{šakutes} \textsuperscript{(12)}, \VUgras{peilius}
\textsuperscript{(13)}, \VUgras{dubenis} \textsuperscript{(15)}, --- kurių nebereikia.

%%K t04-tit-3 sub
\VUpk{10.2pt}{40}{Giminė.}
\VUsaut{3.0mm}

%%K t04-c1-08-1 p
8. --- Visi atrodo linksmi, o šypsena, su kuria jaunikio \VUgras{motina}
\textsuperscript{(17)} žiūri į savo vaikus, rodo jos laimę.

%%K t04-c1-09-1 p
9. --- Tos vestuvės sujungia dvi turtingas krašto šeimas. Jonas, \VUgras{jaunikis}
\textsuperscript{(18)}, yra didelio žemvaldžio \VUgras{sūnus,} ir jis tampa sumanaus
pramonininko \VUgras{žentu.}

%%K t04-c1-10-1 p
10. --- \VUgras{Sutuoktiniai} jauni, ir vis dėlto jie seniai buvo
\VUgras{susižadėję.} \VUgras{Jaunoji moteris} \textsuperscript{(19)}, Morta, žavinga
savo \VUgras{balta suknele,} papuošta apelsinmedžio žiedais.

%%K t04-c1-11-1 p
11. --- Jos \VUgras{uošvis} \textsuperscript{(20)} labai laimingas turėdamas tokią
malonią \VUgras{marčią,} o Jono \VUgras{anyta} \textsuperscript{(21)} šypsosi matydama
tokią gražią savo \VUgras{dukterį.}

%%K t04-c1-12-1 p
12. --- Stalo gale sėdi kiti \VUgras{svečiai} \textsuperscript{(22)}:
\VUgras{giminės, draugai, pusbroliai, pusseserės.} \VUgras{Vyriausiasis padavėjas}
\textsuperscript{(23)}, su servetėle po pažastimi, uoliai jiems patarnauja.

%%K t04-tit-4 sub
\VUpk{10.2pt}{40}{Draugai.}
\VUsaut{3.0mm}

%%K t04-c1-13-1 p
13. --- Dešinėje stalo pusėje štai \VUgras{panelė} \textsuperscript{(24)}, jaunosios
\VUgras{draugė,} paskui jos \VUgras{brolis,} kuris yra jos \VUgras{garbės palydovas}
\textsuperscript{(25)}. Jis šnekučiuojasi su savo \VUgras{garbės panele}
\textsuperscript{(26)}, jaunikio seserimi, kuri per tas vestuves tampa Mortos
\VUgras{svaine.} Ji žavinga jaunuolė. Ji turi gražių \VUgras{papuošalų, perlų
vėrinį} ir auksinę \VUgras{apyrankę.}

%%K t04-c1-14-1 p
14. --- Šalia jų ponas, kuris stovi ir kelia savo taurę, yra šeimos
\VUgras{draugas} \textsuperscript{(27)}. Jis vienas iš \VUgras{jaunikio liudytojų.}
Jis \VUgras{sako tostą,} geria už jaunavedžių \VUgras{sveikatą} ir linki jiems ilgos
laimės. Jis apsirengęs iškilminga apranga, fraku su \VUgras{lakuotais batais}
\textsuperscript{(28)}. Jis palydi \VUgras{senelę} \textsuperscript{(29)}, kuri yra
\VUgras{našlė.}

%%K t04-c1-15-1 p
15. --- Jaunuolis su \VUgras{fraku} \textsuperscript{(31)} ir plonais juodais ūsais
yra Jono \VUgras{svainis} \textsuperscript{(30)}. Jis įžymus inžinierius. Jis sėdi
šalia labai elegantiškos \VUgras{jaunos ponios} \textsuperscript{(33)}, Mortos
\VUgras{vyresniosios sesers} ir motinos to dailaus mergyčios, kuri ateina, duodama
ranką savo pusbroliui, paduoti gėlės ir \VUgras{palinkėti} jam \VUgras{labo vakaro.}
Ta ponia turi labai gražius \VUgras{auskarus} ir elegantišką \VUgras{galvos
apdangalą.}

%%K t04-c1-16-1 p
16. --- Mažasis pusbrolis, toks keistas su savo \VUgras{palaidine}
\textsuperscript{(36)} ir pūstomis \VUgras{kelnaitėmis} \textsuperscript{(37)}, labai
apgailėtinas. Jis \VUgras{našlaitis} \textsuperscript{(35)}. Labai jaunas jis neteko
tėvo ir motinos. Jo dėdė yra jo \VUgras{globėjas} \textsuperscript{(38)} ir liko
viengungis, kad galėtų auklėti vargšą vaiką. Jis geras žmogus. Jis kiek plikas ir
labai trumparegis; jis nešioja \VUgras{(nosies) pensnė.}

%%K t04-tit-5 sub
\VUsaut{2.4mm}
\VUpk{10.2pt}{40}{Desertas.}
\VUsaut{3.0mm}

%%K t04-c1-17-1 p
17. --- Už pokylininkų, ant stalo, uždengto \VUgras{staltiese,} paruoštas
\VUgras{desertas} \textsuperscript{(39)}. Dubenyje štai vaisiai: \VUgras{persikai}
\textsuperscript{(40)}, \VUgras{kriaušės} \textsuperscript{(41)}, \VUgras{obuoliai}
\textsuperscript{(42)}, \VUgras{abrikosai} \textsuperscript{(43)} ir
\VUgras{vynuogės} \textsuperscript{(44)}. Šalia \VUgras{ananasai}
\textsuperscript{(45)}, gražus \VUgras{pyragas} \textsuperscript{(46)},
\VUgras{likerio buteliai} \textsuperscript{(48)} ir \VUgras{šampanas}
\textsuperscript{(49)}, taip pat švarių \VUgras{lėkščių} \textsuperscript{(47)}
stulpeliai.

%%K t04-c1-18-1 p
18. --- \VUgras{Rūsininkas} \textsuperscript{(50)} atkemša seno vyno
\VUgras{butelį} \textsuperscript{(52)} \VUgras{kamščiatraukiu} \textsuperscript{(51)}.
\VUgras{Kamštis} \textsuperscript{(53)} atrodo labai sunkiai ištraukiamas. Tas
rūsininkas turi po pažastimi \VUgras{servetėlę} \textsuperscript{(54)}.

%%K t04-c1-19-1 p
19. --- Po vakarienės ponai eis į rūkomąjį, o ponios eis rengtis baliui.

%%K t04-tit-6 sub
\VUsaut{5.0mm}
\VUcentre{11.4pt}{0}{\textit{Antroji scena.}}
\VUsaut{1.6mm}
\VUpk{11.4pt}{40}{Senelio gimtadienis.}
\VUsaut{2.6mm}

%%K t04-tit-7 sub
\VUpk{10.2pt}{40}{Apsilankymas.}
\VUsaut{3.0mm}

%%K t04-c2-01-1 p
1. --- Praėjo dešimt metų, Jono tėvai paseno ir labai myli savo tris vaikaičius: du
\VUgras{anūkus} \textsuperscript{(2)} ir vieną \VUgras{anūkę} \textsuperscript{(3)}.
Kadangi šiandien yra \VUgras{senelio gimtadienis,} vaikai ateina palinkėti jiems
gero vardadienio.

%%K t04-c2-02-1 p
2. --- Mažasis Petras dar labai jaunas, jam tik treji metai. Jis mato artinantis
\VUgras{katiną} \textsuperscript{(1)}. Jis nori paglostyti jo gražų šilkinį
\VUgras{kailį} ir ilgą \VUgras{uodegą;} jis nepagalvoja, kad po grakščiomis
\VUgras{letenėlėmis} slepiasi pikti smaili nagai, kurie galbūt jį subraižys.

%%K t04-c2-03-1 p
3. --- Jo sesuo Gabrielė, kuri duoda jam ranką, daug rimtesnė, o Marcelis su savo
dideliu \VUgras{apsiaustu} \textsuperscript{(4)} ir odiniu \VUgras{diržu}
\textsuperscript{(5)} atrodo kiek suvyrėjęs, nepaisant trumpų kelnaičių, kurios
leidžia matyti jo \VUgras{kojines} \textsuperscript{(6)}.

%%K t04-c2-04-1 p
4. --- Jų tėvas apsivilko savo storą žieminį \VUgras{paltą} \textsuperscript{(7)} su
\VUgras{kailine apykakle} \textsuperscript{(8)}, o savo \VUgras{kišenėje}
\textsuperscript{(9)} jis turi skarelę. Jo žmona turi savo veltinę skrybėlę,
papuoštą \VUgras{plunksnomis} \textsuperscript{(10)}, savo \VUgras{švarkelį}
\textsuperscript{(11)}, savo \VUgras{kailinį boa} \textsuperscript{(12)} ir savo
\VUgras{movą} \textsuperscript{(13)}.

%%K t04-c2-05-1 p
5. --- Labai šalta, nes viešpatauja žiema. \VUgras{Sniegas} \textsuperscript{(19)}
krito visą dieną, ir namų stogai visai balti. Kartu su \VUgras{naktimi} šaltis
darosi dar aštresnis. Danguje \VUgras{mėnulis} \textsuperscript{(17)} ir
\VUgras{žvaigždės} \textsuperscript{(18)} šviečia ir perveria \VUgras{tamsą.}

%%K t04-tit-8 sub
\VUsaut{4.2mm}
\VUpk{10.2pt}{40}{Liga.}
\VUsaut{3.0mm}

%%K t04-c2-06-1 p
6. --- Vargšas \VUgras{aklasis} \textsuperscript{(20)}, kuris eina per aikštę
priešais \VUgras{vaistinę} \textsuperscript{(16)}, vedamas savo \VUgras{pudelio}
\textsuperscript{(21)}, labai nelaimingas. Justina, \VUgras{tarnaitė}
\textsuperscript{(14)}, neša iš virtuvės \VUgras{puodelį} \textsuperscript{(15)}
labai karštos \VUgras{arbatžolių nuoviro,} dosniai pasaldinto.

%%K t04-c2-07-1 p
7. --- \VUgras{Senelė} \textsuperscript{(23)}, kuri nori pasiimti nuo stalo savo
\VUgras{akinius} \textsuperscript{(26)}, kad perskaitytų \VUgras{receptą}
\textsuperscript{(27)}, kurį \VUgras{gydytojas} \textsuperscript{(25)} ką tik
parašė, kyla iš savo krėslo, atsiremdama į savo \VUgras{lazdą}
\textsuperscript{(24)}.

%%K t04-c2-08-1 p
8. --- Ji dažnai kenčia nuo \VUgras{negalavimų, svaigulių, slogų, dantų skausmų;}
jos kojos silpnos, o akys nebėra labai geros. Nepaisant to, ji noriai pašiepia savo
seną \VUgras{gydytoją,} kuris visada nori jai išrašyti \VUgras{mikstūrą}
\textsuperscript{(28)}. Šiandien ji labai džiaugsminga, nes turi aplink save savo
jauną šeimą.

%%K t04-c2-09-1 p
9. --- Gerai uždarytame kambaryje jauku. Marmuriniame \VUgras{židinyje} \textsuperscript{(29)}
liepsnoja \VUgras{graži ugnis} \textsuperscript{(30)}, ir jautiesi laimingas švelnioje
\VUgras{šviesoje} \VUgras{lempos} \textsuperscript{(31)}, kurią ką tik uždegė.

%%K t04-tit-9 sub
\VUsaut{4.2mm}
\VUpk{10.2pt}{40}{Senelis.}
\VUsaut{3.0mm}

%%K t04-c2-10-1 p
10. --- Net \VUgras{senelis} \textsuperscript{(33)} pamiršo, kad sirgo, kad jo
\VUgras{reumatas} laiko jį prirakintą prie \VUgras{krėslo} \textsuperscript{(34)} ir
kad dar vakar jį kamavo \VUgras{karščiavimas ir alpimai.} Jis nebeprisimena, kad per
paskutinę žiemą sirgo \VUgras{plaučių uždegimu} ir kad kimus ir varginantis
\VUgras{kosulys} labai jį kankino.

%%K t04-c2-10-2 p
Ir vis dėlto labai sunkiai jis buvo išgydytas. Nuo šiol jam kiek geriau. Bet koja,
kurią jis remia į \VUgras{pagalvę} \textsuperscript{(38)}, padėtą ant mažo
\VUgras{suoliuko} \textsuperscript{(39)}, irgi jį skauda.

%%K t04-c2-11-1 p
11. --- Kai jis rūpestingai suvyniotas į savo šiltą \VUgras{chalatą}
\textsuperscript{(35)} su storomis perlamutrinėmis \VUgras{sagomis}
\textsuperscript{(36)} ir kai turi šalia
savęs, kad jam paskaitytų laikraštį, mažąją \VUgras{našlaitę} \textsuperscript{(40)},
apsivilkusią balta \VUgras{kykele,} mėlyna \VUgras{suknele} \textsuperscript{(42)} ir
\VUgras{pelerina} \textsuperscript{(41)}, jis nesiskundžia.

%%K t04-c2-12-1 p
12. --- Kasmet, kai \VUgras{kalendorius} \textsuperscript{(43)} vėl rodo pirmosios
\VUgras{gruodžio} dienos \VUgras{datą,} jis nekantriai laukia savo \VUgras{anūkų,}
nes žino, kad jie nepamirš tos svarbios \VUgras{datos.}

%%K t04-c2-13-1 p
13. --- Jis jaudindamasis prisimena tuos labai senus laikus, kai pats ateidavo
palinkėti geros šventės garsiajam imperijos \VUgras{maršalui,} kurio
\VUgras{portretas} \textsuperscript{(44)} kabo ant sienos ir kuris yra jo vaikų
\VUgras{prosenelis.}

\VUsaut{6.0mm}
\VUfilet{22mm}
